\apartado{Problema, objetivos y alcance}
\subsection{Planteamiento del problema}
Un estudiante puede distribuir su trabajo entre archivos, apuntes, calendarios y herramientas de consulta. Sin una estructura común, resulta difícil identificar qué material pertenece a una materia, qué pendientes están próximos y qué información se comparte. La incorporación de IA añade una dificultad: una respuesta convincente puede no corresponder a los apuntes o exponer contenido que debía permanecer privado.

El problema de ingeniería consiste en construir una base que organice esa información, conserve su pertenencia y permita ampliar funciones sin comprometer identidad, persistencia o trazabilidad. Desde la asignatura, interesa explicar cómo los servicios cloud resuelven necesidades concretas y qué responsabilidades conserva el equipo.

\subsection{Objetivo general}
Diseñar e iniciar una plataforma universitaria de estudio mediante una arquitectura cloud administrada, estableciendo una base reproducible y un espacio académico privado que puedan validarse de manera incremental.

\subsection{Objetivos específicos}
\begin{enumerate}
\item Definir contratos comunes de datos, errores y autorización para interfaz y backend.
\item Evaluar el aislamiento de recursos interactivos y la extracción documental, sin confundir prototipos con módulos terminados.
\item Construir el acceso del piloto y la organización de materias, temas, periodos y agenda.
\item Relacionar identidad, almacenamiento, contenedores, escalamiento y observabilidad con los requisitos.
\item Establecer evidencias reproducibles y controles de consumo para decidir la continuidad de cada fase.
\end{enumerate}

\subsection{Delimitación}
La \textbf{fase 1 corresponde a H0 y T01--T04}; la \textbf{fase 2 corresponde a H1 y T05--T08}. No se reinterpretan como H1 y H2. El alcance incluye desarrollo local y preparación arquitectónica para cloud. Aprovisionar recursos, publicar servicios o consumir IA pagada requiere las condiciones operativas correspondientes y evidencia adicional.

La tutoría completa, carga documental de extremo a extremo, generación y guardado de recursos, práctica evaluada, progreso, pasaporte y funciones sociales pertenecen a hitos posteriores. Tampoco se incluyen pagos, comunidad pública, integración institucional ni ejecución arbitraria de código en servidor. El nombre definitivo y el catálogo de estampillas siguen pendientes.
